\section{Related Work}
PET reconstruction is commonly formulated from a Poisson likelihood, with expectation--maximization providing a seminal estimator \citep{shepp1982maximum}. Modern iterative reconstruction incorporates physical corrections and priors \citep{qi2006iterative}. Post-reconstruction denoising remains attractive when raw projection data are unavailable, but it cannot recover information absent from the acquisition.

The Anscombe transform maps a Poisson variable to an approximately unit-variance Gaussian variable \citep{anscombe1948transformation}. More accurate inverse formulae reduce low-count bias \citep{makitalo2011optimal}; nevertheless, the conventional algebraic inverse is retained here to match the defined experimental model. TV regularization formalizes edge-preserving smoothing \citep{rudin1992nonlinear}, while higher-order models mitigate piecewise-constant staircasing \citep{bredies2010total}. Spatial adaptation can interpolate between these behaviors using evidence from a pilot image.

Projected gradient descent enforces simple convex constraints \citep{bertsekas1999nonlinear}. Barzilai--Borwein steps approximate curvature from successive iterates at negligible memory cost \citep{barzilai1988two}, while Armijo backtracking provides a descent safeguard \citep{nocedal2006numerical}. Learned PET methods have also shown promise \citep{gong2019pet}, but variational baselines retain advantages in interpretability, reproducibility, and limited-data settings.
